\section{The controlled base potential}\label{sec:base-potential}

We now sum the logarithmic tests to prove
Proposition~\ref{base:existence}.  Fix parameters and an integer $Q$
for which Theorem~\ref{density:main} and Lemma~\ref{tests:later}
hold, and write $k_j=Q^j$.  For a finite amplitude $L>0$, set
\begin{equation}\label{base:potential}
 \phi(z)=\psi(z)+L\sum_{j=1}^\infty U_{k_j}(z).
\end{equation}
We first justify this definition.  We then choose $L$ to exclude
the required pluriharmonic minorant, and finally prove the uniform
chart estimates needed by the metric construction.

\begin{lemma}[Smooth convergence]\label{base:convergence}
For every finite $L>0$, the series in~\eqref{base:potential}
converges locally in $C^\infty(\mathbb B)$.  Its sum is smooth,
satisfies $\phi(0)=0$, and has complex Hessian $h\ge I$.
\end{lemma}

\begin{proof}
Fix $0<r_0<r_1<1$.  Around each point of $|z|\le r_0$, the complex
polydisk of radius $(r_1-r_0)/(2\sqrt2)$ lies in $|z|<r_1$.
Cauchy estimates there, together with Lemma~\ref{peaks:bounds}, give
\[
 \sup_{|z|\le r_0}|D^mP_k(z)|
 \le C_{m,r_0,r_1}r_1^k
\]
for every fixed derivative order $m$.  Here and below $D^m$ denotes
any real partial derivative of total order $m$.
The values of $P_k$ remain in a fixed bounded disk.
Every positive-order derivative of $F_\epsilon(P_k)$ is a sum of
products containing at least one positive-order derivative of
$P_k$ or its conjugate.  The bounded derivatives of $F_\epsilon$
therefore give the same exponential estimate for $D^mU_k$.
For $m=0$, use~\eqref{tests:lipschitz}.
Since $\sum_j r_1^{Q^j}<\infty$, the series converges in every local
$C^m$ norm.

Each summand vanishes at zero and is plurisubharmonic by
\eqref{tests:hessian}.  Termwise differentiation is justified by
the local $C^2$ convergence, and hence
\[
 \partial\bar\partial\phi
 =I+L\sum_j\partial\bar\partial U_{k_j}\ge I.
\]
\end{proof}

\begin{lemma}[Absence of a pluriharmonic minorant]\label{base:minorant}
If
\begin{equation}\label{base:amplitude}
 L>\frac{5\log Q}{d_3},
\end{equation}
there are no $H\in\mathcal O(\mathbb B)$ and $C_H\in\mathbb R$
satisfying~\eqref{base:no-minorant}.
\end{lemma}

\begin{proof}
Put
\begin{equation}\label{base:radii}
 r_N=1-\frac{\eta}{k_N}.
\end{equation}
For every $j\le N$, Bernoulli's inequality gives
\[
 r_N^{k_j}
 =\left(1-\frac{\eta}{k_N}\right)^{k_j}
 \ge1-\eta\frac{k_j}{k_N}\ge1-\eta.
\]
Thus all of the first $N$ modes satisfy the hypothesis of
\eqref{tests:later-negative} at this same radius.
For $j=N+m>N$, on the other hand,
\[
 r_N^{k_j}\le e^{-\eta k_j/k_N}=e^{-\eta Q^m}.
\]
By the uniform polynomial bound and~\eqref{tests:lipschitz},
\begin{equation}\label{base:tail}
 \sum_{j>N}\sup_{\xi\in S}|U_{k_j}(r_N\xi)|
 \le C_\epsilon\sum_{m\ge1}e^{-\eta Q^m}
 =:C_{\rm tail}<\infty,
\end{equation}
independently of $N$.  Because $W_N$ is a positive probability,
the same bound controls the weighted tail, without requiring a
uniform pointwise bound on $W_N$.

The base potential contributes only one logarithm:
\[
 \psi(r_N\xi)=-\log(1-r_N^2)
 \le\log\frac1{1-r_N}=\log(k_N/\eta).
\]
Combining this with~\eqref{tests:later-negative} and
\eqref{base:tail} yields
\begin{align}
 &\int_S\left(\phi(r_N\xi)
                  +4\log\frac1{1-r_N}\right)W_N\,d\sigma \notag\\
 &\hspace{1cm}\le
 5\log(k_N/\eta)+L(-d_3N+C_{\rm tail}) \notag\\
 &\hspace{1cm}=
 N(5\log Q-Ld_3)+5\log(1/\eta)+LC_{\rm tail}
 \longrightarrow-\infty.
 \label{base:negative-means}
\end{align}

By Theorem~\ref{density:main}, the measures
$\mu_N=W_N\,d\sigma$ are positive representing probabilities.
Apply Lemma~\ref{peaks:mean-criterion} to these measures, the radii
$r_N$, and the continuous function
$u(z)=\phi(z)+4\log(1/(1-|z|))$.
The limit~\eqref{base:negative-means} is exactly its hypothesis,
so the required minorant cannot exist.
\end{proof}

It remains to control the metric jets as the chart centers
approach the boundary.  The value of $\phi$ need not be bounded.
For its derivatives, the low degrees become almost constant in a
centered chart, while the high degrees are exponentially small.

\begin{lemma}[Uniform chart estimates]\label{base:chart-estimates}
For every fixed finite $L>0$, the potential~\eqref{base:potential}
has uniformly bounded positive-order derivatives through order four
at the origins of all centered ball charts.  Consequently its
complex Hessian satisfies~\eqref{base:control} and
\eqref{base:jets}, with $c=1$ and a finite constant $C$.
\end{lemma}

\begin{proof}
After a unitary change of coordinates, put the center at $(r,0)$,
where $0\le r<1$, and let $\delta=1-r$.  Rotating $P_k$ preserves its
degree and its uniform homogeneous bound.  Use the chart $T_r$ of
\eqref{base:charts}, on the fixed complex polydisk
\[
 D_\rho=\{|\zeta_1|<\rho,\ |\zeta_2|<\rho\},
 \qquad \rho=\frac1{16}.
\]
The denominators in $T_r$ have modulus at least $1-\rho$.
We estimate the polynomial in the two ranges of $k\delta$.

\paragraph{High degrees: $k\delta>1$.}
The ball automorphism identity gives, on $D_\rho$,
\[
 1-|T_r(\zeta)|^2
 =\frac{(1-r^2)(1-|\zeta|^2)}{|1+r\zeta_1|^2}
 \ge \frac{1-2\rho^2}{(1+\rho)^2}\,\delta.
\]
It follows that $1-|T_r(\zeta)|\ge c\delta$, with a fixed $c>0$.
Thus
\[
 |P_k(T_r(\zeta))|\le C_P e^{-ck\delta}.
\]
Cauchy estimates in this fixed polydisk bound every fixed-order
holomorphic derivative at zero, including order zero, by
$C_m e^{-ck\delta}$.

\paragraph{Low degrees: $k\delta\le1$.}
Consider the polydisk in the original coordinates
\[
 A_k=\left\{|z_1-r|<\frac1k,\quad
                 |z_2|<\frac1{\sqrt k}\right\}.
\]
It may extend outside the ball, but $P_k$ is an entire homogeneous
polynomial.  The bound of Lemma~\ref{peaks:bounds} therefore holds
there as well.  In fact
\[
 |z|^2\le(1+1/k)^2+1/k\le1+4/k,\qquad
 |P_k(z)|\le C_P(1+2/k)^k\le C_Pe^2.
\]
Cauchy estimates on the concentric half-polydisk give
\[
 |\partial_1P_k|\le Ck,\qquad
 |\partial_2P_k|\le C\sqrt k.
\]
Meanwhile the chart displacements satisfy
\[
 |T_r^1(\zeta)-r|
 \le\frac{2\rho}{1-\rho}\delta,\qquad
 |T_r^2(\zeta)|
 \le\frac{\sqrt2\rho}{1-\rho}\sqrt\delta.
\]
Both numerical coefficients are less than $1/2$.  As
$\delta\le1/k$, the chart image lies in the half-polydisk just
used.  Integrating the gradient along the straight segment from
$(r,0)$ to $T_r(\zeta)$ gives
\[
 |P_k(T_r(\zeta))-P_k(r,0)|
 \le C(k\delta+\sqrt{k\delta})
 \le2C\sqrt{k\delta}.
\]
This segment stays in the half-polydisk by convexity.  Apply
Cauchy estimates on $D_\rho$ to the holomorphic difference
$P_k\circ T_r-P_k(r,0)$.  Every positive-order derivative at zero
is bounded by $C_m\sqrt{k\delta}$.  We make no such assertion for
the zeroth-order value in this range.

\paragraph{Composition and summation.}
In both ranges the values of $P_k\circ T_r$ stay in a fixed bounded
disk.  Each positive-order real derivative of
$U_k\circ T_r=F_\epsilon(P_k\circ T_r)$ is a sum of products
containing differentiated factors of $P_k\circ T_r$ or its
conjugate.  For $1\le m\le4$, the preceding bounds and the smooth
composition estimate give
\begin{equation}\label{base:scale-jets}
 |D^m(U_k\circ T_r)(0)|
 \le C_{\epsilon,m}
 \begin{cases}
  \sqrt{k\delta},&k\delta\le1,\\
  e^{-ck\delta},&k\delta>1.
 \end{cases}
\end{equation}
Indeed every term contains at least one differentiated factor,
and a product of several factors obeys the same bound because the
displayed scale factor is at most one.  This proves the estimate
for all real derivatives, hence for all pure and mixed complex
derivatives of these orders.

Let $J$ be the largest index with $Q^J\delta\le1$, or put $J=0$
if there is no such index.  The low-degree contribution is bounded
by the geometric sum
\[
 \sum_{j\le J}\sqrt{Q^j\delta}
 \le\frac1{1-Q^{-1/2}}.
\]
For the remaining degrees, the first value of $Q^j\delta$ exceeds
one, so
\[
 \sum_{j>J}e^{-cQ^j\delta}
 \le\sum_{m\ge0}e^{-cQ^m}<\infty.
\]
Both bounds are uniform in the chart center.  Termwise
differentiation is legitimate by Lemma~\ref{base:convergence}.
Multiplication by the fixed amplitude $L$ preserves these bounds.

Finally, the base potential transforms as
\[
 \psi\circ T_r
 =-\log(1-|\zeta|^2)-\log(1-r^2)
      +\log|1+r\zeta_1|^2.
\]
Its positive-order derivatives through order four at zero are
uniformly bounded.  The last term is pluriharmonic, so the complex
Hessian of $\psi\circ T_r$ is exactly the standard ball Hessian in
these coordinates.  The derivatives of $\phi\circ T_r$ needed for
the pulled-back Hessian components and their first and second
derivatives have orders two, three and four.  The bounds just
proved, together with $h\ge I$ from
Lemma~\ref{base:convergence}, establish the assertion.
\end{proof}

\begin{proof}[Proof of Proposition~\ref{base:existence}]
Choose $a,\alpha,s_1,f,b,R_0$ as in Lemma~\ref{peaks:profiles}.
Proposition~\ref{tests:negative} first chooses $\epsilon,\eta$ and
then permits any sufficiently large spacing $D_0$.  Fix such a
$D_0$ and its degree threshold $k_*$, and fix
$d_3=d_2/2>0$.  The density theorem holds for every $Q$ above a
threshold depending only on the fixed patch data and spacing.
Its derivative constant $K_*$ is fixed before that threshold is
chosen.  We may therefore choose one integer $Q$ satisfying the
density threshold, $Q\ge k_*$, and the Fourier-error threshold of
Lemma~\ref{tests:later}.

Now fix a finite amplitude $L$ satisfying~\eqref{base:amplitude}.
Lemma~\ref{base:convergence} constructs the smooth potential with
$h\ge I$; Lemma~\ref{base:minorant} proves the required obstruction;
and Lemma~\ref{base:chart-estimates} supplies its finite uniform
chart bounds.  The chart constant is allowed to depend on these
fixed choices, including $L$.  The metric parameter $\lambda$ in
Proposition~\ref{metric:transfer} is chosen only afterward from
these bounds, so none of the choices above depends on $\lambda$.
\end{proof}
